\documentclass[runningheads]{llncs}

\usepackage[T1]{fontenc}

% T1 fonts will be used to generate the final print and online PDFs,
% so please use T1 fonts in your manuscript whenever possible.
% Other font encondings may result in incorrect characters.

\usepackage[utf8]{inputenc}
\usepackage{CJKutf8}

\hyphenation{co-domain}


\title{Finite Functional Programming}
\subtitle{or, LAMBDA: The Ultimate Predicate}
%\title{Finite Functional Logic Programming \\via Pointed Types}
% If the paper title is too long for the running head, you can set
% an abbreviated paper title here
%\titlerunning{Finite Functional Logic Programming via Pointed Types}

\author{Michael Arntzenius\orcidID{0009-0002-0417-5636}
  \and Max Willsey\orcidID{0000-0001-8066-4218}
}
%
\authorrunning{M. Arntzenius}
% First names are abbreviated in the running head.
% If there are more than two authors, 'et al.' is used.
%
\institute{University of California, Berkeley, USA\\
\email{daekharel@gmail.com}}


%% ---- Packages ----
\usepackage{color}
\usepackage[dvipsnames,svgnames]{xcolor}

%\usepackage{anyfontsize} % silence font size not available warnings by autoscaling fonts
\usepackage{amssymb}            % \multimap, \mathbb
%\usepackage{amsmath}
%\usepackage{latexsym,stmaryrd}
%\usepackage{mathtools}          % \dblcolon
\usepackage{mathpartir}         % mathpar
\usepackage{nccmath}            % fleqn (for code environment)

%\usepackage{minted}
\usepackage{listings}           % lstlisting
\lstset{language=prolog, columns=fullflexible,
  basicstyle=\ttfamily, commentstyle=\color{Green}}

\usepackage{url,hyperref}
% If you use the hyperref package, please uncomment the following two lines
% to display URLs in blue roman font according to Springer's eBook style:
\renewcommand\UrlFont{\color{blue}\rmfamily}
\urlstyle{rm}

\usepackage[nameinlink,noabbrev]{cleveref}
\newcommand\crefrangeconjunction{--} % use en-dashes for ranges.


%% TEMPORARY: different fonts because I hate reading computer modern
\usepackage[scaled=0.97]{XCharter}\linespread{1.0333}
%\usepackage[scaled=0.965,proportional,semibold]{sourcesanspro}
\usepackage[osf,scaled=1.0333]{AlegreyaSans}
%\usepackage[scaled=1.08]{biolinum}
\usepackage[scaled=1.021]{inconsolata}
\usepackage[xcharter,scaled=1.0556,]{newtxmath}
%\usepackage{mathastext}


%% ---------- Code environment ----------
\usepackage{environ}            %\NewEnviron
\newlength\codeoffset
\setlength\codeoffset{0pt}

\NewEnviron{code}{
  \begin{fleqn}[\codeoffset]
    \begin{flalign*}
      \setlength\arraycolsep{0pt} % necessary for flush left alignment.
      \begin{array}{l}\BODY\end{array}
  \end{flalign*}
  \end{fleqn}
}


%% ---------- \negphantom command ----------
%% https://tex.stackexchange.com/questions/316426/negative-phantom-inside-equations
\makeatletter
\newlength{\negph@wd}
\DeclareRobustCommand{\negphantom}[1]{%
  \ifmmode
    \mathpalette\negph@math{#1}%
  \else
    \negph@do{#1}%
  \fi
}
\newcommand{\negph@math}[2]{\negph@do{$\m@th#1#2$}}
\newcommand{\negph@do}[1]{%
  \settowidth{\negph@wd}{#1}%
  \hspace*{-\negph@wd}%
}\makeatother


%% ---- Commands ----
\newcommand\strong[1]{{\bfseries#1}}
\newcommand\todo[1]{{\color{IndianRed}#1}}
\newcommand\TODO[1]{\todo{\strong{TODO}: #1}}
\newcommand\xxx{\todo{\strong{XXX}}}

\newcommand\fslang{\ensuremath{\lambda_{\mathrm{FS}}}}
\newcommand\catname[1]{\ensuremath{\mathrm{#1}}}
\newcommand\Point{\ensuremath{\catname{Set}_{*}}}
\DeclareMathOperator{\support}{supp}

\newcommand\den[1]{\llbracket{#1}\rrbracket}
\newcommand\bnfeq{\Coloneqq}
\newcommand\bnfor{\mathrel{\,|\,}}

\newcommand\name[1]{\ensuremath{\mathit{#1}}}
%\newcommand\<\;                 %function application space
\newcommand\<\:                 %function application space
\newcommand\fnspace{\<}         %space after a binder
\newcommand\fn[1]{\lambda{#1}.\fnspace}
\newcommand\fa[1]{\forall{#1}.\fnspace}
\newcommand\ex[1]{\exists{#1}.\fnspace}

%\newcommand\cto{\Rightarrow}    % "cartesian \to"
\newcommand\cto\to
\newcommand\lto\multimap        % "linear \to"
\newcommand\fto\Rightarrow
%\renewcommand\fto{\underset{\textsf{fin}}{\to}}           % finite map
%\newcommand\fto{\rightarrow}    % "finitely supported \to"
\newcommand\x\times
\newcommand\with{\mathbin{\&}}
\newcommand\ox\otimes
\newcommand\typename\text
\newcommand\bool{\name{bool}}
\newcommand\maybe[1]{\ensuremath{\typename{maybe}~{#1}}}

%% need \usepackage{accent} for these
%\newcommand\undertype[2]{\underaccent{\textsf{\upshape\scshape#1}}{#2}}
%\newcommand\eqt[1]{\undertype{eq}{#1}}

\newcommand\G\Gamma
\newcommand\D\Delta
\renewcommand\O\Omega
\renewcommand\L\Lambda

%\newcommand\emptycx\varepsilon
\newcommand\emptycx{{\cdot}}
\newcommand\of{\mathbin{:}}
\newcommand\cxsep{\mathbin{/}}
\newcommand\hyp[2]{{#1} \of {#2}}

\newcommand\cJ[3]{{#3} \,\vdash\, {#1} : {#2}}
\newcommand\J[5]{{#3} \cxsep {#4} \cxsep {#5} \,\vdash\, {#1} : {#2}}
%\newcommand\pJ[5]{{#3} \cxsep {#4} \cxsep {#1} : {#2} \,\vdash\, {#5}}

\newcommand\kword[1]{\ensuremath{\mathrm{#1}}}
\newcommand\const\name
%\newcommand\cname\text          %constructor name
%\newcommand\cname\name
\newcommand\cname\textsf
%\newcommand\cnil{\cname{nil}}
\newcommand\cjust{\cname{just}}
%\newcommand\ctrue{\cname{true}}
%\newcommand\cfalse{\cname{false}}

\newcommand\nil{\name{nil}}
\newcommand\true{\name{true}}
\newcommand\false{\name{false}}

\newcommand\letin[1]{\kword{let}~{#1}~\kword{in}~}
\newcommand\leteq[2]{\letin{{#1} = {#2}}}
\newcommand\tjust[1]{\cjust\<{#1}}
\newcommand\letjust[3]{\leteq{\tjust {#1}}{#2}{#3}}
\newcommand\letpair[4]{\leteq{({#1}, {#2})}{#3}{#4}}
\newcommand\tabsurd{\kword{absurd}\<}
\newcommand\tand{\mathrel{\kword{and}}}
\newcommand\tor{\mathrel{\kword{or}}}
%% \renewcommand\tand{~\kword{and}~}
%% \renewcommand\tor{~\kword{or}~}
%\newcommand\twhen[2]{\kword{when} \<({#1}) \<{#2}}
\newcommand\twhen[2]{\kword{if}~ {#1} ~\kword{then}~ {#2}}
%\newcommand\twhen[2]{{#1} \Rightarrow {#2}}
%\newcommand\twhen[2]{{#2} \mathrel{\kword{when}} {#1}}
%\newcommand\teq[2]{\name{eq} \<{#1} \<{#2}}
\newcommand\teq[2]{{#1} = {#2}}
%\newcommand\teq[2]{{#1} \mathrel{\kword{is}} {#2}}
%\newcommand\teq[2]{{#2} \gets {#1}}
\newcommand\tcomment[1]{\color{black!60}\textsf{\#\ #1}}

%% \renewcommand\tand{\mathrel{\mathrm{\&\&}}}
%% \renewcommand\tor{\mathrel{\mathrm{||}}}
%% \newcommand\tand{\mathrel{\mathbf{and}}}
%% \newcommand\tor{\mathrel{\mathbf{or}}}

\newcommand\pwild\textunderscore
\newcommand\nat{\mathbb{N}}


\begin{document}
\maketitle              % typeset the header of the contribution

\begin{abstract}
We propose unifying functional and logic/relational programming by treating predicates/relations as functions equipped with their support: the set of inputs whose output is nonzero.
For instance, the logic programming language Datalog is a language of finitely supported boolean functions.
Finite support allows representing the function, not as code, but as data---as a table of input-output pairs---making bulk operations more efficient.
Aggregation and weighted logic programming are naturally handled by generalizing the definition of support from booleans to other pointed sets.
We refer to the combination of finitely supported functions, represented as data, with higher order functions, represented as code, as \emph{finite functional programming.}
To this end, we demonstrate a simple type system to check finite support, using a combination of graded coeffects to check variable grounding and relevance types to model pointed sets.

%% %% ---------- OLD ABSTRACT ----------
%%   I propose unifying functional and logic/relational programming by treating predicates/relations as functions equipped with their support---the set of inputs whose output is non-false.
%%   For instance, the logic programming language Datalog, in which all relations are finite, is a language of finitely supported boolean functions.
%%   %
%%   Finite support allows representing the function, not as code, but as data---as a table of input-output pairs---making bulk operations more efficient.
%% %
%%   If we combine finitely supported functions, represented as data, and higher order functions, represented as code, in the same language, we get what I call \emph{finite functional programming.}
%%   %
%%   To this end, I demonstrate a simple type system to check finite support \todo{THIS IS NOT QUITE RIGHT, IT'S THE STATIC GRADING NOT THE RUNTIME EFFECT THAT MATTERS: by treating variable grounding as a graded effect.}
%%   %Aggregation and other Datalog extensions (such as bag and semiring semantics) are naturally handled by generalizing from boolean functions to arbitrary pointed output types.
%%   %% \todo{do I actually do this? $\rightarrow$}
%%   %% I sketch future work on semantics and implementation strategies, particularly for recursion.

\keywords{Logic programming \and Functional programming \and Categorical semantics \and Linear types \and Coeffects}
\end{abstract}


\section{Logic programming is (not) functional programming}


%\section{Support and pointed sets}
%\section{What is the point of support?}
%\section{What is support? What is the point?}
\section{Support, pointedly}


\section{Why not just...?}


\section{A finitely supported map from examples to booleans}
%\section{\fslang\ as she is wrote}
%\section{A finite set of finite functional programs}
%\section{Some finite functional programs}


%\section{Pointed sets and relevance}
%\section{Points, relevantly}
%\section{The relevance of relevance}
%\section{The importance of being relevant}
\section{The relevance of being relevant}

\section{Context hopping}


%% \section{Programming with pointed sets}
%\section{\fslang}
%\section{\fslang\ as she is typed}

%\section{Decidably pointed sets}
%\section{Recovering maybes}
\section{Curing maybe-addiction with class}


%\section{Semantics of \fslang}
\section{What does it all mean?}

\end{document}
